% Source: https://artofproblemsolving.com/community/c3h3134177_2012_lmt_spring_guts_round_p1p12__lexington_mathematical_tournament
% Stable identity: aops-topic-3134177-q07
The crease is the perpendicular bisector of $AP$, so the marked point is its midpoint. The map $P\mapsto\dfrac{A+P}{2}$ is a homothety with center $A$ and factor $\dfrac{1}{2}$. The locus is that circle with the point $A$ omitted, because $P=A$ was excluded. Its closure is the full circle with one-half the original radius, whose enclosed area is $\dfrac{1}{4}\cdot80=\boxed{20}$.
\begin{asy}
usepackage("amsmath");
size(7cm);
pair A=(2,0),P=2dir(120),M=(A+P)/2;draw(circle((0,0),2));draw(circle((1,0),1));draw(A--P);dot(M);label("$A$",A,E);label("$P$",P,NW);label("$M$",M,N);
\end{asy}
